\section{Methods}
\subsection{Population and episode anchor}
Polish-speaking adults aged at least 18 who actively directed AI-generated music within the past 30 days are eligible, including abandoned or disappointing sessions. Passive listening alone is insufficient. Select the most recent eligible episode before presenting items, irrespective of affective intent. A session begins with intentional initiation of a generative music task and ends when the user stops that task. Recruit via permission-based online channels with neutral language; no claims of emotional improvement. The convenience sample does not estimate population prevalence.

\subsection{Stage A: content development}
Item construction follows staged development principles \parencite{boateng2018}, drawing on relevance, comprehensiveness and comprehensibility criteria from COSMIN \parencite{terwee2018} as adapted to this nonclinical instrument. The 35 original Polish items (five per provisional domain) are candidate content only. A versioned alternative pool proposes 12 revisions while retaining the original 35 identifiers and source wording. The changes distinguish symbolic expression, mood-function wording, repeated attention and within-episode avoidance; they are not approved content. Qualitative themes from P02 and literature inform revision; no themes are fabricated. Maintain a traceability register linking item, domain, rationale, qualitative evidence when available and revision decision.

Recruit 6--10 experts spanning psychometrics, music psychology, emotion regulation and human--AI interaction, and conduct two rounds of 6--10 cognitive interviews including novice and unsuccessful creators. Experts rate relevance and clarity separately on four-point scales; report item-level proportions rating 3 or 4, denominators and exact binomial intervals. Values below .78 prompt review, not automatic deletion. Interviews use paraphrasing, retrieval and response-mapping probes; attend to overlap between distraction and avoidance and between reflection and repetition. Maintain a content-coverage matrix; no item survives only because of expert agreement.

\begin{figure}[htbp]
\centering
\begin{tikzpicture}[
  stage/.style={draw, rounded corners=2pt, align=center, minimum width=33mm,
    minimum height=13mm, font=\small},
  arrow/.style={-{Latex[length=2mm]}}]
\node[stage] (content) {P02 evidence, theory\\and candidate items};
\node[stage, right=4mm of content] (review) {Expert review and\\cognitive interviews};
\node[stage, right=4mm of review] (efa) {Feasibility pilot,\\then ordinal EFA};
\node[stage, right=4mm of efa] (cfa) {Frozen model,\\independent CFA};
\draw[arrow] (content) -- (review);
\draw[arrow] (review) -- (efa);
\draw[arrow] (efa) -- (cfa);
\end{tikzpicture}
\caption{Prospective development sequence. The record linking each item to qualitative findings also documents contradictory evidence and gaps in domain coverage. P02 evidence is pending.}
\end{figure}

Expert ratings and participant interpretations address complementary questions.
An item can receive high relevance ratings while respondents consistently
recall the initial intention or evaluate the quality of the track. In that case,
revise the wording and examine the interpretation again in the second interview
round. Link each revision to an item version and to the qualitative evidence
that motivated the change.

\subsection{Stage B: pilot and exploratory study}
A feasibility pilot of 30--50 adults tests comprehension, response distributions, burden and skip patterns. Pilot records are not pooled with factor analysis samples. The provisional EFA target is 400 eligible records (maximum 480 submissions or eight weeks). Freeze the item pool before collection; no outcome-based stopping. Investigate sparsity, missingness and floor effects; low endorsement alone is not sufficient to delete meaningful content. Record episode completion and days since episode; an abandoned session is not equivalent to zero on all functions.

\subsection{Stage C: independent confirmation}
After EFA, freeze item wording, domain assignment, estimator, missing-data handling, factor model, comparators and primary validity contrasts before recruiting an independent CFA sample (provisional target 500 eligible records, maximum 600 submissions or eight weeks). Use non-identifying study codes and self-reported prior participation to prevent overlap across P05 samples; no IP retention. P02/P03/P04 data may inform development but are not independent validation evidence. Recruitment targets are contingent on simulation and practical feasibility.

\subsection{Response format and external measures}
For the selected episode, rate the extent to which the activity served each specified function from 0 (not at all) to 4 (very much); provide a separate cannot-assess/prefer-not-to-answer option coded missing. Items apply even without a finished output. No reverse wording solely for balancing. Comparator selection: authorised Polish music mood regulation measure (MMR/B-MMR where suitable), authorised rumination and avoidance measures, and separate initial motive and recalled clarity indicators as in P03. HUMS and ERP-IA-6 are candidate benchmarks only if age, language, permission and construct suitability are established. Do not reproduce protected items without authorisation. An exploratory incremental-validity comparison with a suitable general AI regulation scale is contingent on these prerequisites. A subsample target of 100 CFA participants will rerate the same episode after 7--14 days; interpret agreement as stability of recollection, not trait stability.

\subsection{Ethics and open science}
Obtain institutional ethics determination and informed consent before recruitment. Participants may omit any item or withdraw before submitting anonymous responses. Do not request lyrics, prompts, audio, trauma narratives, handles or clinical diagnoses. Explain that AI-MRFS is not diagnostic or therapeutic. Store research records in approved institutional storage; only approved aggregate outputs, synthetic data and non-sensitive analysis code belong in Git. For optional retest, store contacts separately with restricted access and remove them after follow-up under the approved retention policy. Record provenance and do not describe multiple analyses from one dataset as independent replication. Preregister separately before EFA and CFA and report all deviations. No approval or participant recruitment is claimed.
